\documentclass[a4paper]{article} %

\date{} %

% Please, do not change any of the above lines

\usepackage{amsmath,amssymb,hyperref} % add other LaTeX packages you need here.

\begin{document}
\setcounter{page}{1} % Please, do not delete this line

% Your title
\title{Residuated partially ordered algebras,\\ logic and preclones}

% Authors
\author{Peter Jipsen \\ %
       Chapman University, Orange, California, USA\\ \\ % Affiliation 1
       \texttt{jipsen@chapman.edu} % Only one corresponding e-mail, please
       }%

\maketitle

% The abstract


%\subsection*{Abstract}
\noindent
In the first lecture we introduce partially ordered algebras with operations that are monotone or antitone in each argument. The standard examples are algebras of binary relations, residuated po-semigroups and residuated lattices, as well as many po-subvarieties and their connection with substructural logic and algebraic proof theory. 
The second installment concentrates on the structure theory of residuated po-algebras, considering poset decompositions and po-Plonka sums. 
In the last lecture we focus on the general theory of S-preclones for a finite monoid S, which generalizes clone theory to the po-algebra setting and beyond.

In more detail, we aim to cover the following topics:

\medskip

\noindent
\textbf{Lecture 1: Tonicities, inequational logic, and partially ordered varieties} (based on standard sources and joint work with Marta Bilkova and Melissa Sugimoto).
We begin with the universal algebraic theory of partially ordered algebras ($\tau$-po-algebras) specified by a set $\mathcal F$ of function symbols and a tonicity $\tau: \mathcal F \to \{+, -\}^*$. This framework allows primitive operations to be monotone ($+$) or antitone ($-$) at specified argument positions, naturally capturing structures such as left-residuated po-monoids, residuated po-semigroups, residuated lattices, po-groups, tense posets, and algebras of binary relations. In this setting, inequational logic is treated as a 2-dimensional deductive system within Abstract Algebraic Logic (AAL). We consider $\tau$-preorders (the order analogues of congruences) and cover order-theoretic counterparts of the fundamental isomorphism and correspondence theorems, Birkhoff's subdirect representation theorem, and the order $\mathsf{H}\text{-}\mathsf{S}\text{-}\mathsf{P}$ and $\mathsf{S}\text{-}\mathsf{L}_{\to}\text{-}\mathsf{P}$ characterization theorems for $\tau$-povarieties and $\tau$-quasivarieties. We also discuss aspects of algebraizability of substructural logics and algebraic proof theory.

\medskip

\noindent
\textbf{Lecture 2: Structure theory via po-P\l onka sums of residuated po-semigroups} (based on joint work with Stefano Bonzio, Jos\'e Gil-F\'erez, Adam P\v renosil, Melissa Sugimoto and Giuseppe Zecchini).
The second lecture focuses on the structure theory of residuated partially ordered semigroups $\langle A, \leqslant, \cdot, \backslash, /\rangle$, where $x \cdot y \leqslant z \iff x \leqslant z/y \iff y \leqslant x \backslash z$. We investigate the role of positive central idempotents $p \in \mathrm{ZE}^+ \mathbf{A}$ ($a \leqslant pa = ap$ and $p^2 = p$) and balanced residuated semigroups (satisfying $x \backslash x \approx x / x =: 1_x$ and $y \leqslant 1_x y$). To reconstruct such algebras from their fibers $\mathbf{A}_p = \{a \in A : 1_a = p\}$, we generalize the classical P\l onka sum construction to \emph{semilattice-directed systems of metamorphisms} $\Xi = \{\xi_{pq}: \mathbf{A}_p \looparrowright \mathbf{A}_q : p \preccurlyeq q \in I\}$. Unlike homomorphisms, metamorphisms handle each operation and relation coordinatewise via distinct families of maps $\langle \varphi_{pq}, \psi_{pq}\rangle$. We establish compatibility conditions under which the reconstructed partial order
\[
a \leqslant b \iff \varphi_{pr}(a) \leqslant_r \psi_{qr}(b) \quad (\text{where } r = p \vee q)
\]
yields a valid po-algebra. This establishes that every \emph{steady} residuated semigroup decomposes into a po-P\l onka sum of integrally closed residuated monoids (and into Brouwerian semilattices in the idempotent commutative case).

\medskip

\noindent
\textbf{Lecture 3: $S$-preclones, $S$-relations, and the ${}^S\mathrm{Pol}\text{--}{}^S\mathrm{Inv}$ Galois connection} (based on joint work with Erkko Lehtonen and Reinhard P\"oschel).
In the final lecture, we present the general theory of $S$-preclones for an arbitrary finite monoid $S$ of signa, generalizing classical clone theory to signed operations. An $S$-operation $f: A^n \to A$ carries an argument signum $\mathrm{sgn}(f) = (s_1, \dots, s_n) \in S^n$. $S$-preclones are sets of $S$-operations containing the identity $\mathrm{id}_A$ and closed under cyclic shifts ($\zeta$), argument transpositions ($\tau$), addition of fictitious signed arguments ($\nabla^s$), identification of equal-signum arguments ($\Delta$), and composition ($\circ$). Dually, an $m$-ary $S$-relation is a family $\boldsymbol{\varrho} = (\varrho_s)_{s \in S}$ with $\varrho_s \subseteq A^m$. We introduce the $S$-preservation relation
\[
f \rhd^S \boldsymbol{\varrho} \iff \forall s \in S : f(\varrho_{s_1 s}, \dots, \varrho_{s_n s}) \subseteq \varrho_s
\]
and study the induced Galois connection ${}^S\mathrm{Pol}\text{--}{}^S\mathrm{Inv}$. For finite $A$, we prove the Galois closure theorems: ${}^S\langle F \rangle = {}^S\mathrm{Pol}\,{}^S\mathrm{Inv}\,F$ and ${}^S[Q] = {}^S\mathrm{Inv}\,{}^S\mathrm{Pol}\,Q$. Furthermore, we analyze the lattice ${}^S\mathcal{L}_A$ of all $S$-preclones on $A$, showing that it is atomic and coatomic with finitely many atoms and coatoms, explore embeddings of the classical clone lattice $\mathcal{L}_A$, and highlight the motivating case of $\pm$-preclones over the group $S = \{+, -\}$ characterizing monotone and antitone operations on posets.


\begin{thebibliography}{99}

\bibitem{BGJPS}
S.~Bonzio, J.~Gil-F\'erez, P.~Jipsen, A.~P\v{r}enosil, and M.~Sugimoto.
\newblock Balanced residuated partially ordered semigroups.
\newblock \emph{preprint} \url{https://arxiv.org/abs/2505.12024}, 2025.

\bibitem{JLP2024}
P.~Jipsen, E.~Lehtonen, R.~P\"oschel.
\newblock S-preclones and the Galois connection $^S$Pol--$^S$Inv, Part I.
\newblock \emph{Algebra Universalis}, 85:34, 2024. \url{https://doi.org/10.1007/s00012-024-00863-7}

\bibitem{JTV2021}
P.~Jipsen, O.~Tuyt, and D.~Valota.
\newblock The structure of finite commutative idempotent involutive residuated lattices.
\newblock \emph{Algebra Universalis}, 82(4):57, 2021. 
\url{https://digitalcommons.chapman.edu/scs_articles/736/}

\bibitem{Pi04}
D.~Pigozzi.
\newblock Partially ordered varieties and quasivarieties.
\newblock Technical report, Iowa State University, 2004. 
\small\url{https://www.researchgate.net/publication/250226836_partially_ordered_varieties_and_quasivarieties}

\end{thebibliography}

\end{document}
